% New post-CP16 work. Do not integrate until complete independent audits.
\subsection{Geometric generic nonordinary principally polarized surfaces}
\label{subsec:generic-nonordinary-ppav-groups-candidate}
\status{Candidate: primary source statements and proofs opened; complete new assembly awaits independent audit. These are geometric generic statements and finite-degree open consequences, not assertions for every fixed nonordinary surface.}

\begin{lemma}[Frame component counts under nonproper specialization]
\label{lem:frame-components-specialization-candidate}
Let $B$ be an excellent normal integral scheme, let
$U\subset\mathbf P_B(\mathcal W)$ be an open whose geometric
fibres are nonempty, and let $\mathcal V$ be a rank-$r$
$\mathbf F_2$ local system on $U$. Write $G_{\bar b}$ and
$G_{\bar\eta}$ for its geometric images on a geometric fibre
and on the geometric generic fibre. Then
\[
                         |G_{\bar\eta}|\ge |G_{\bar b}|.
\]
If every geometric-fibre image is, up to conjugation, contained
in a fixed subgroup $P\subset\GL_r(\mathbf F_2)$ and one
geometric fibre has image $P$, the geometric generic image is
$P$. If that equality holds generically, it holds on every
geometric fibre over a nonempty open of $B$.
\end{lemma}

\begin{proof}
Use the full frame cover $Z\to U$, of degree
$|\GL_r(\mathbf F_2)|$. Its geometric fibre over $b$ has
$|\GL_r|/|G_{\bar b}|$ connected components, because the
image acts freely on frames. Choose a finite separable
extension of $k(B)$ defining all components of
$Z_{\bar\eta}$, and normalize $B$ in it. The normalization
$B'\to B$ is finite surjective by excellence. The pullback
$U'$ is a normal integral open of a projective bundle over
$B'$. The finite etale cover $Z'$ is normal, so closures of
its generic components are disjoint open-and-closed
components. Each surjects onto $U'$, since its image is both
open and closed. They all meet every geometric fibre of
$U'\to B'$. Hence the number of geometric generic components
is at most the number on any geometric special fibre,
giving the asserted inequality.

For the last assertion, shrink $B$ so that the normalization
is finite etale, and apply
\cite[Lemma~37.27.5]{Stacks0559} to each component of $Z'$.
Its generic fibre is geometrically irreducible; it is
geometrically normal because it is finite etale over an open
of projective space. The lemma gives an open of $B'$ on
which all these fibres are geometrically irreducible.
Remove from $B$ the finite image of the complementary
closed subset. The remaining open is nonempty and all its
geometric frame fibres have exactly the generic number of
components. Their images therefore have the generic order.
Together with the prescribed subgroup bound, this proves
equality up to conjugation. This argument uses full frames;
it requires no choice of invariant-covector basis over $B$.
\end{proof}

\begin{lemma}[Constant-rank relative unit systems in ppav families]
\label{lem:ppav-family-constant-rank-Lang-candidate}
Let $A\to B$ be a family of principally polarized abelian
surfaces in characteristic two, supplied with a relatively
ample line bundle $M$ representing the polarization, and
suppose every geometric fibre has the same two-rank
$f\in\{0,1\}$. Fix $d\ge1$ and put $b=9d^2$.
On the open $U_d$ of smooth curves in $|M^{3d}|$ having
maximal relative stable rank, the relative Lang system is a
finite etale local system of rank
\[
 r=\begin{cases}b-1,&d\text{ odd},\\ b-3+f,&d\text{ even}.
 \end{cases}
\]
Every geometric fibre of $U_d\to B$ is a nonempty integral
open of a projective space. Its geometric image is contained
in $\GL_r(\mathbf F_2)$ in odd degree and in the pointwise
stabilizer of $f+1$ independent covectors in even degree.
\end{lemma}

\begin{proof}
Cohomology and base change for the ample line bundle gives
a rank-$b$ vector bundle of sections, and the universal
hypersurface sequence gives a vector bundle of relative
coherent cohomology of rank $b-1$. The rank formula and
nonemptiness of the smooth maximal-rank open on every
geometric fibre are Theorem~\ref{thm:DE-A3-ppav-ranks}.
All smaller stable ranks are detected by minors of a fixed
high Frobenius iterate: an iterate of exponent at least
$b-1$ has image equal to the stable image on each geometric
coherent fibre.

On $U_d$ the linearization of that iterate has constant rank
$r$, so its image is a vector subbundle, using the usual
invertible-minor presentations. Frobenius restricts to an
isomorphism of its Frobenius pullback with itself, by the
same fibrewise rank statement and Nakayama's lemma. Its
Lang fixed scheme is therefore finite etale and gives the
claimed local system. This constructs a single system over
the family even when the nilpotence index of the complementary
part changes. It avoids identifying the first image with
the stable image.

The last subgroup bounds are the entire invariant-covector
theorem, Theorem~\ref{thm:nonordinary-ppav-all-covectors},
applied to each geometric fibre. In even degree its
dimension is one for $f=0$ and two for $f=1$.
\end{proof}

\begin{lemma}[The needed genus-two strata and their specializations]
\label{lem:ppav-nonordinary-strata-candidate}
Over $\overline{\mathbf F}_2$, the two-rank-one stratum of
principally polarized abelian surfaces is smooth and
irreducible. It contains the mixed product locus.
Every geometric irreducible component of the two-rank-zero
stratum has a normal one-dimensional parameter scheme
dominating it and containing a point representing the
product principally polarized surface $E_0^2$.
\end{lemma}

\begin{proof}
The primary source is
\cite[Corollary~8.4, printed pp.~607--608; Theorem~11.5 and
Example~11.6, p.~620; Example~12.2, pp.~623--624]{EkedahlVanderGeer2009Strata}.
The four genus-two final strata are, in order, the ordinary
stratum, the two-rank-one stratum, the two-rank-zero
$a$-number-one stratum, and the superspecial stratum.
Their final flag strata map finite etale of degree one onto
them. Corollary~8.4 gives smoothness of the open flag strata
and normality of their closures. Theorem~11.5 and
Example~11.6 give irreducibility of the two-rank-at-most-one
locus, hence of its dense two-rank-one open. These statements
hold in characteristic two: the large-characteristic
condition in Theorem~11.4(ii) concerns the converse
component-counting assertion, not Theorem~11.5.

For the two-rank-zero stratum, take a component of the
closed one-dimensional flag stratum $\overline U_{s_2}$.
It is normal by Corollary~8.4 and its dense open maps
degree one to the corresponding $a$-number-one component.
Each such component meets the zero-dimensional superspecial
flag stratum by
\cite[Proposition~6.1 and its proof, printed p.~601]{EkedahlVanderGeer2009Strata}.
No irreducibility of the full supersingular locus is asserted
or needed. The flag bundle over the abelian moduli space is
projective, and the closed stratum is closed in that bundle.
The forgetful map on a chosen normal curve component is
therefore proper. It is generically degree one onto its
image by Example~12.2. A nonconstant map from an integral
curve has zero-dimensional fibres, so this proper map is
quasi-finite and hence finite. It is the normalization of
its image. In particular the superspecial flag point just
obtained is an actual point of this normal parameter curve;
we have not replaced the intersecting ppav component locus
itself by a smooth scheme without a map.

We verify that the underlying principally polarized surface
at such a point is $E_0^2$ in characteristic two. By
Lemma~\ref{lem:ppav-surface-dichotomy}, an indecomposable
principally polarized surface has a smooth genus-two theta
curve and is its Jacobian. A smooth genus-two curve cannot
have Cartier identically zero. Indeed in its function field
the kernel of Cartier on rational differentials is
one-dimensional over the square subfield. If both independent
canonical differentials were killed, their ratio would be a
square. The degree-two canonical map to $\mathbf P^1$
would then be purely inseparable. For a one-variable
function field $K$ over a perfect algebraically closed
field, $[K:K^2]=2$; a purely inseparable degree-two
canonical map would force $K^2=k(\mathbf P^1)$ and hence
genus zero, a contradiction.
Thus an $a$-number-two surface is decomposable. Both
elliptic factors have $a$-number one, hence are supersingular.
Over an algebraically closed field of characteristic two
there is a unique supersingular elliptic isomorphism class,
represented by $E_0:y^2+y=x^3$. This last assertion follows
also by normalizing a Weierstrass equation with $a_1=0$
and $a_3\ne0$ and making successive translations of $x,y$.
The polarization in the decomposable case is the product
polarization, proving the claimed specialization.

All these arguments may be made on a fine prime-to-two
level-five moduli scheme. For two-rank one, choose a
connected etale chart through the desired mixed product;
its smooth component maps openly, hence dominantly, to
the irreducible stratum. We do not assume that a level
cover is irreducible. For two-rank zero, pull back the
normal flag component to that finite etale level cover
and take a component. It still dominates the original
component and meets a point above the superspecial
specialization, since a finite etale component over a
normal connected base is surjective. A local etale
neighbourhood supplies a line bundle representing the
principal polarization: the relative Picard torsor is
smooth and the line bundle can be rigidified at zero.
No symmetric-theta torsor is required in characteristic two.
These changes preserve normality and generic dominance.
\end{proof}

\begin{theorem}[Generic full images on nonordinary ppav surfaces]
\label{thm:generic-nonordinary-ppav-groups-candidate}
Let $(A,M)$ be a geometric generic principally polarized
abelian surface on a two-rank-$f$ stratum, $f\in\{0,1\}$,
in characteristic two. For $f=0$ this means the geometric
generic point of any irreducible component.
For every $d\ge1$, the relative system on the maximal-rank
smooth open of the complete linear system $|M^{3d}|$ has
\[
 G_{\mathrm{geom},d}=
 \begin{cases}
 \GL_{9d^2-1}(\mathbf F_2),&d\text{ odd},\\[2pt]
 \operatorname{Hom}(\mathbf F_2^{f+1},\mathbf F_2^{9d^2-4})
       \rtimes\GL_{9d^2-4}(\mathbf F_2),&d\text{ even}.
 \end{cases}
\]
For any perfect field of definition with such geometric
generic moduli, the arithmetic image equals the geometric
image. For the embedding by $M^3$, the same assertion for
ambient degree-$d$ hypersurfaces holds whenever ambient
restriction is surjective, in particular for
$d\ge\max(1,\operatorname{reg}(\mathcal I_A)-1)$.
\end{theorem}

\begin{proof}
For $f=1$, base change the stratum to an algebraically
closed field containing a transcendental ordinary elliptic
modulus. Choose its mixed product point and a normal
etale chart through it as in the preceding lemma.
Every degree at that point has the asserted full group
by Theorem~\ref{thm:generic-mixed-product-all-groups-candidate}.
The relative ranks on the entire stratum are the same as
at this point. The family Lang lemma and the full frame
component-count lemma therefore give at least that group
order at the geometric generic point. The entire
invariant-covector theorem gives precisely that full
stabilizer as an upper bound. The two bounds coincide.
Dominance of the chart identifies this geometric generic
group with that of the original moduli stratum; an
algebraically closed extension does not change the image
of a finite etale local system.

For $f=0$, use the normal parameter component and its
superspecial product point supplied by the preceding
lemma. The full rank is constant at every geometric
point, including this specialization. The product groups
in every degree are Theorem~\ref{thm:ss-product-all-full-groups}.
The same frame-component inequality and invariant-covector
upper bound prove the assertion at the generic point of
each component. This use of a normalization is precisely
why singularities or multiple branches of a supersingular
coarse moduli locus cause no gap.

In both cases the entire geometric invariant space is
arithmetic over a perfect field, by
Theorem~\ref{thm:nonordinary-ppav-all-covectors}.
Its pointwise stabilizer is the displayed geometric image,
so arithmetic equality follows. The ambient restriction
claim uses only the stated surjectivity and the integral
fibres of a linear projection; it does not assume
characteristic-two projective normality of the cubic
theta embedding.
\end{proof}

\begin{corollary}[Open and finite-field consequences for finitely many degrees]
\label{cor:generic-nonordinary-ppav-finite-degrees-candidate}
Fix a finite set of positive degrees and one of the strata
or components in the preceding theorem. On a nonempty
Zariski open of that stratum or component, every geometric
surface has the displayed full groups in all the chosen
degrees. Consequently there are examples defined over
finite extensions of $\mathbf F_2$ with those groups in
all the chosen complete linear systems.
No claim that one finite-field point works for every
degree follows from this corollary.
\end{corollary}

\begin{proof}
For a single degree, apply the last assertion of the
frame-component lemma to the universal family on a
normal chart. After a finite separable extension defining
its geometric generic frame components, shrink to make
that cover finite etale and make every component's fibres
geometrically irreducible using Stacks~\cite{Stacks0559}.
The resulting full-image property descends to an open
of the original chart and then to an open of the moduli
component: an etale map is open, and the property depends
only on the geometric polarized surface. For a normalized
two-rank-zero curve, first restrict to the dense normal
open of its image, where the normalization is an
isomorphism. Intersect the finitely many resulting dense
opens for the prescribed degrees.

The moduli charts, universal families, finite covers and
their finitely many open conditions descend from
$\overline{\mathbf F}_2$ to a finite field. A nonempty
finite-type scheme over a finite field has a closed
point with finite residue field. Its polarized surface
gives the asserted examples, after a further finite
extension if needed to choose a representing line bundle.
Only finitely many opens are intersected; no countable
intersection or reversed specialization is used.
\end{proof}
